\documentclass[../main.tex]{subfiles}

\begin{document}

\ifSubfilesClassLoaded{

    %\tableofcontents
    
    %\setcounter{chapter}{}
}{}

\chapter{ Root Systems }
% 代靖涵 25120222201319
\section{Definition of Root Systems}

\begin{definition}{\(  s_v  \) }{}
    Let \(E\) be a finite-dimensional real vector space endowed with an inner product written \((-, -)\). Given a non-zero vector \(v \in E\), let \(s_v\) be the reflection in the hyperplane normal to \(v\). Thus \(s_v\) sends \(v\) to \(-v\) and fixes all elements \(y\) such that \((y, v) = 0\). 
\end{definition}

\begin{proposition}{}{}

    \begin{itemize}
        \item 
\[
s_{v}(x)=x-\frac{2(x,v)}{(v,v)}v\quad\text{for all }x\in E
\]
\item \(s_{v}\) preserves the inner product, that is,

\[
(s_{v}(x),s_{v}(y))=(x,y)\quad\text{for all }x,y\in E.
\]
\item  As a useful convention , we write\[
\langle x,v\rangle:=\frac{2(x,v)}{(v,v)},
\] then \[
s_{v}= x-\left<x,v \right>v
\]


    \end{itemize}
    

\end{proposition}

\begin{remark}
    The symbol \(  \left<x,v \right>  \) is only linear with respect to the first variable, \(  x  \).  
\end{remark}

\begin{definition}{Root system}{}
A subset $R$ of a real vector space $E$ is a \dfntxt{root system} if it satisfies the following axioms.
\begin{itemize}
  \item[(R1)] $R$ is finite, it spans $E$, and it does not contain $0$.
  \item[(R2)] If $\alpha \in R$, then the only scalar multiples of $\alpha$ in $R$ are $\pm \alpha$.
  \item[(R3)] If $\alpha \in R$, then the reflection $s_{\alpha}$ permutes the elements of $R$.
  \item[(R4)] If $\alpha, \beta \in R$, then $\langle \beta, \alpha \rangle \in \mathbf{Z}$.
\end{itemize}
The elements of $R$ are called \dfntxt{roots}.
\end{definition}

\section{First Steps in the Classification}
\begin{lemma}{Finiteness Lemma}{}
Suppose that \(R\) is a root system in the real inner-product space \(E\). Let \(\alpha,\beta\in R\) with \(\beta\neq\pm\alpha\). Then
\[
\langle\alpha,\beta\rangle\langle\beta,\alpha\rangle\in\{0,1,2,3\}.
\]
\end{lemma}

\begin{proofsketch}
    将 \(  \left<\alpha ,\eta  \right>\left<\beta ,\alpha  \right>  \)展开写成内积, 
\end{proofsketch}
\section{Basis for Root Systems}

\begin{definition}{}{}
A subset \(B\) of \(R\) is a \dfntxt{base} for the root system \(R\) if
\begin{enumerate}
  \item[(B1)] \(B\) is a vector space basis for \(E\), and
  \item[(B2)] every \(\beta \in R\) can be written as \(\beta = \sum_{\alpha \in B} k_{\alpha}\alpha\) with \(k_{\alpha} \in \mathbf{Z}\), where all the non-zero coefficients \(k_{\alpha}\) have the same sign.
\end{enumerate}
\end{definition}
\begin{remark}
    We say that a root \(\beta \in R\) is \dfntxt{positive with respect to \(B\)} if the coefficients given in (B2) are positive, and otherwise it is \dfntxt{negative with respect to \(B\)}.
\end{remark}
\end{document}